%% start of file `template.tex'.
%% Copyright 2006-2015 Xavier Danaux (xdanaux@gmail.com).
%
% This work may be distributed and/or modified under the
% conditions of the LaTeX Project Public License version 1.3c,
% available at http://www.latex-project.org/lppl/.


\documentclass[11pt,letterpaper,roman,colorlinks=true,urlcolor=brown]{moderncv}

% moderncv themes
\moderncvstyle{banking}
\moderncvcolor{blue}
%\renewcommand{\familydefault}{\sfdefault}
%\nopagenumbers{}

% page geometry
\usepackage[scale=0.75]{geometry}

% personal data
\name{Yuxiao}{Qu}
\homepage{cohenqu.github.io}
\extrainfo{Google Scholar: \url{https://scholar.google.com/citations?user=X5SFuyYAAAAJ\&hl=en}}

% bibliography style
\renewcommand*{\bibliographyitemlabel}{[\arabic{enumiv}]}
%----------------------------------------------------------------------------------
%            content
%----------------------------------------------------------------------------------
\begin{document}
%\begin{CJK*}{UTF8}{gbsn}                          % to typeset your resume in Chinese using CJK
%-----       resume       ---------------------------------------------------------
\makecvtitle

\vspace{-1.4cm}

% \section{Research Interests}
% My research focuses on the intersection of reinforcement learning and foundation models. I am passionate about developing innovative algorithms and scalable systems to tackle diverse decision-making challenges in complex reasoning and agentic environments. I am especially interested in test-time compute optimization, generalization through abstraction, and combining supervised and reinforcement learning for large language models.

\section{Education}
\cventry{May 2024 - Present}{PhD in Machine Learning}{Carnegie Mellon University}{Pittsburgh, PA}{4.22/4.33}{Advisor: \href{https://aviralkumar2907.github.io/}{Aviral Kumar} \& \href{https://www.cs.cmu.edu/~rsalakhu/}{Ruslan Salakhutdinov}}
\vspace{0.2cm}
\cventry{Aug 2023 - Present}{Master of Science in Machine Learning }{Carnegie Mellon University}{Pittsburgh, PA}{4.28/4.33}{Advisor: \href{https://aviralkumar2907.github.io/}{Aviral Kumar}}
\vspace{0.2cm}
\cventry{Sep 2020 – Dec 2022}{Bachelor of Science in Computer Science and Mathematics}{University of Wisconsin - Madison}{Madison, WI}{4.00/4.00}{}
\vspace{0.2cm}
\cventry{Sep 2017 – May 2020}{Major: Computer Science}{The Chinese University of Hong Kong}{Sha Tin, Hong Kong}{}{}

\section{Working Experiences}
\cventry{Apr 2023 – Aug 2023}{Research Software Engineer}{Morgridge Institute for Research}{Madison, WI}{}{}
\cventry{Nov 2025 – Present}{Research Intern}{NVIDIA}{Remote}{}{}

\section{Awards and Fellowships}
\begin{itemize}
    \item 2025 Amazon PhD Fellowship In AI \hfill 2025
    \item 2025 Qualcomm Innovation Fellowship (\textit{Finalist}) \hfill 2025
    \item 2021 ICPC North Central NA Regional Contest (\textit{Rank 19th}) \hfill 2021
    \item Professor Charles K. Kao Research Exchange Scholarship \hfill 2021
    \item Best Project Award in Undergraduate Summer Research Internship \hfill 2019
\end{itemize}


\section{Publications}
(* indicates Equal Contribution)
\begin{enumerate}
    \item \textbf{Learning to Reason on Hard Problems with Privileged On-Policy Exploration}\\
    \textbf{Yuxiao Qu}*, Amrith Setlur*, Virginia Smith, Ruslan Salakhutdinov, Aviral Kumar\\
    \textcolor{red}{\textit{Oral}} at 5th MATH-AI Workshop, NeurIPS 2025
    \item \textbf{CaRT: Teaching LLM Agents to Know When They Know Enough}\\
    Grace Liu*, \textbf{Yuxiao Qu}*, Jeff Schneider, Aarti Singh, Aviral Kumar
    \item \textbf{Learning to Discover Abstractions for LLM Reasoning}\\
    \textbf{Yuxiao Qu}*, Anikait Singh*, Yoonho Lee*, Amrith Setlur, Ruslan Salakhutdinov, Chelsea Finn, Aviral Kumar\\
    \textcolor{red}{\textit{Spotlight}} at Workshop on Efficient Systems for Foundation Models, ICML 2025\\
    \textcolor{red}{\textit{Oral}} at RAM 2: Reasoning, Attention \& Memory – 10 Years On, COLM 2025
    \item \textbf{Optimizing Test-Time Compute via Meta Reinforcement Fine-Tuning}\\
    \textbf{Yuxiao Qu}*, Matthew Y.R. Yang*, Amrith Setlur, Lewis Tunstall, Edward Emanuel Beeching, Ruslan Salakhutdinov, Aviral Kumar
    \\
    \textcolor{red}{\textit{Oral}} at Workshop on Foundation Models in the Wild, ICLR 2025\\
    International Conference on Machine Learning (ICML), 2025
    \item \textbf{Harnessing Webpage UIs for Text-Rich Visual Understanding}\\
    Junpeng Liu, Tianyue Ou*, Yifan Song*, \textbf{Yuxiao Qu}*, Wai Lam, Chenyan Xiong, Wenhu Chen, Graham Neubig, Xiang Yue
    \\
    International Conference on Learning Representations(ICLR), 2025
    \item \textbf{Recursive Introspection: Teaching Language Model Agents How to Self-Improve}\\
    \textbf{Yuxiao Qu}, Tianjun Zhang, Naman Garg, Aviral Kumar
    \\
    \textcolor{red}{\textit{Oral}} at Workshop on Large Language Models and Cognition, ICML 2024\\
    Advances in Neural Information Processing Systems (NeurIPS), 2024
    \item \textbf{Guided Data Augmentation for Offline Reinforcement Learning and Imitation Learning}\\
    Nicholas E. Corrado, \textbf{Yuxiao Qu}, John U. Balis, Adam Labiosa, Josiah P. Hanna
    \\
    Reinforcement Learning Conference (RLC), 2024
    \item \textbf{FLEE-GNN: A Federated Learning System for Edge-Enhanced Graph Neural Network in Analyzing Geospatial Resilience of Multicommodity Food Flows}\\
    \textbf{Yuxiao Qu}, Jinmeng Rao, Song Gao, Qianheng Zhang, Wei-Lun Chao, Yu Su, Michelle Miller, Alfonso Morales, Patrick Huber\\
    ACM SIGSPATIAL International Workshop on AI for Geographic Knowledge Discovery, 2023
    \item \textbf{Simulation-Acquired Latent Action Spaces for Dynamics Generalization}\\
    Nicholas E. Corrado, \textbf{Yuxiao Qu}, Josiah P. Hanna\\
    Proceedings of Machine Learning Research (PMLR), 2022
    \item \textbf{Polite or Direct? Conversation Design of a Smart Display for Older Adults Based on Politeness Theory}\\
    Yaxin Hu, \textbf{Yuxiao Qu}, Adam Maus, Bilge Mutlu\\
    ACM CHI Conference on Human Factors in Computing Systems (CHI), 2022
    \item \textbf{TreeNet: Deep Point Cloud Embedding for Routing Tree Construction}\\
    Wei Li, \textbf{Yuxiao Qu}, Gengjie Chen, Yuzhe Ma, Bei Yu\\
    \textcolor{red}{\textit{Best Paper Award}} at ACM Asian and South Pacific Design Automation Conference (ASPDAC), 2021
\end{enumerate}

\section{Research Blog Posts}
\begin{enumerate}
    \item \textbf{IsoCompute Playbook: Optimally Scaling Sampling Compute for RL Training of LLMs} \hfill 2026\\
    Zhoujun Cheng*, Yutao Xie*, \textbf{Yuxiao Qu}*, Amrith Setlur*, Shibo Hao, Varad Pimpalkhute, Tongtong Liang, Feng Yao, Hector Liu, Eric Xing, Virginia Smith, Ruslan Salakhutdinov, Zhiting Hu, Taylor Killian, Aviral Kumar\\
    \href{https://compute-optimal-rl-llm-scaling.github.io/}{Blog Post}
    \item \textbf{QED-Nano: Teaching a Tiny Model to Prove Hard Theorems} \hfill 2026\\
    LM-Provers, \textbf{Yuxiao Qu}, Amrith Setlur, Jasper Dekoninck, Edward Beeching, Jia Li, Ian Wu, Lewis Tunstall, Aviral Kumar\\
    \href{https://huggingface.co/spaces/lm-provers/qed-nano-blogpost}{Blog Post}
    \item \textbf{How to Explore to Scale RL Training of LLMs on Hard Problems?} \hfill November 2025\\
    \textbf{Yuxiao Qu}, Amrith Setlur, Virginia Smith, Ruslan Salakhutdinov, Aviral Kumar\\
    \href{https://blog.ml.cmu.edu/2025/11/26/how-to-explore-to-scale-rl-training-of-llms-on-hard-problems/}{ML@CMU}
    \item \textbf{Optimizing LLM Test-Time Compute Involves Solving a Meta-RL Problem} \hfill January 2025\\
    Amrith Setlur, \textbf{Yuxiao Qu}, Matthew Yang, Lunjun Zhang, Virginia Smith, Aviral Kumar\\
    \href{https://blog.ml.cmu.edu/2025/01/08/optimizing-llm-test-time-compute-involves-solving-a-meta-rl-problem/}{ML@CMU}
\end{enumerate}

\section{Teaching Experiences}
% \subsection{Teaching Assistant}
\begin{itemize}
    \item COMP SCI 577: Introduction to Algorithms (\textit{Honor Session}) \hfill Fall 2022
    \item COMP SCI 537: Introduction to Operating Systems \hfill Fall 2022
\end{itemize}
% \subsection{Employment}
% \cventry{Apr 2023 – Aug 2023}{Research Software Engineer}{Morgridge Institute for Research}{Madison, WI}{}{}

% \section{Research Mentoring}
% % \subsection{Undergraduate & Masters Students}
% \begin{itemize}
%     \item Lehong Wu (2025-; Peking University)
% \end{itemize}
\section{Professional Service}
\begin{itemize}
    \item Reviewer for ICML (2024-2025), NeurIPS (2024-2025), AISTATS (2025), ICLR (2025-2026)
\end{itemize}


\end{document}